\onecolumn

\chapter{Auswertung}
\label{ch:Auswertung}
% Liste der genutzer Formeln für die Fehlerrechnung
\section*{Fehlerrechnung}
Für die statistische Auswertung von $n$ Messwerten $x_i$ werden folgende Größen definiert \cite{errorSkript25}:
\begin{align}
    \bar{x} &= \frac{1}{n} \sum_{i=1}^{n} x_i \vphantom{\sqrt{\sum_i^n}^2} && \text{\textcolor{gray}{Arithmetisches Mittel}} \label{eq:arithmetisches_mittel} \\
    \sigma^2 &= \frac{1}{n-1} \sum_{i=1}^{n} (x_i - \bar{x})^2 \vphantom{\sqrt{\sum_i^n}^2} && \text{\textcolor{gray}{Variation}} \label{eq:variation} \\
    \sigma &= \sqrt{\frac{1}{n-1} \sum_{i=1}^{n} (x_i - \bar{x})^2} \vphantom{\sqrt{\sum_i^n}^2} && \text{\textcolor{gray}{Standardabweichung}} \label{eq:standardabweichung} \\
    \Delta \bar{x} &= \frac{\sigma}{\sqrt{n}} = \sqrt{\frac{1}{n(n-1)} \sum_{i=1}^n(\bar x - x_i)^2} \vphantom{\sqrt{\sum_i^n}^2} && \text{\textcolor{gray}{Fehler des Mittelwerts}} \label{eq:fehler_mittelwert} \\
    \Delta f &= \sqrt{\left(\frac{\partial f}{\partial x} \Delta x\right)^2 + \left(\frac{\partial f}{\partial y} \Delta y\right)^2} \vphantom{\sqrt{\sum_i^n}^2} && \text{\textcolor{gray}{Gauß’sches Fehlerfortpflanzungsgesetz für $f(x,y)$}} \label{eq:gauss_fehlfortpflanzung} \\
    \Delta f &= \sqrt{(\Delta x)^2 + (\Delta y)^2} \vphantom{\sqrt{\sum_i^n}^2} && \text{\textcolor{gray}{Fehler für $f = x + y$}} \label{eq:fehler_summe} \\
    \Delta f &= |a| \Delta x \vphantom{\sqrt{\sum_i^n}^2} && \text{\textcolor{gray}{Fehler für $f = ax$}} \label{eq:fehler_proportional} \\
    \frac{\Delta f}{|f|} &= \sqrt{\left(\frac{\Delta x}{x}\right)^2 + \left(\frac{\Delta y}{y}\right)^2} \vphantom{\sqrt{\sum_i^n}^2} && \text{\textcolor{gray}{relativer Fehler für $f = xy$ oder $f = x/y$}} \label{eq:relativer_fehler} \\
    \sigma &= \frac{|a_{lit} - a_{gem}|}{\sqrt{\Delta a_{lit}^2 + \Delta a_{gem}^2}} \vphantom{\sqrt{\sum_i^n}^2} && \text{\textcolor{gray}{Berechnung der signifikanten Abweichung}} \label{eq:signifikante_abweichung}
\end{align}

\twocolumn

\section{Adiabatenkoeffizient nach Clément-Desormes}
\label{sec:A1}

Fehler des Adiabatenkoeffizienten \(\kappa\) werden via \hyperref[eq:gauss_fehlfortpflanzung]{Gauß'scher Fehlerfortpflanzung (\ref*{eq:gauss_fehlfortpflanzung})} berechnet. Die Formel für \(\kappa\) (\ceq{CD_k_final}) lautet:

\eq{errCD_k}{
    \Delta \kappa_{\text{CD, Luft}} = \sqrt{\frac{\Delta {h_1}^2 \, {h_3}^2 + \Delta{h_3}^2 \, {h_1}^2}{(h_1 - h_3)^4}} \, .
}

\begin{table}[h]
    \centering
    \caption{Ergebnisse der Adiabatenkoeffizienten $\kappa$ der einzelnen Messungen.}
    \label{tab:kappaA1}
    \begin{tabular}{c c}
        \toprule
        Messung & $\kappa_\text{Luft}$ \\
        \midrule
        1 & $1.34 \pm 0.06$ \\
        2 & $1.38 \pm 0.07$ \\
        3 & $1.39 \pm 0.06$ \\
        4 & $1.36 \pm 0.06$ \\
        5 & $1.39 \pm 0.07$ \\
        6 & $1.39 \pm 0.07$ \\
        7 & $1.36 \pm 0.06$ \\
        \bottomrule
    \end{tabular}
\end{table}


Da wir mehrere Messwerte haben, berechnen wir den Mittelwert \(\bar{\kappa}_{\text{CD, Luft}}\) und den Fehler des Mittelwerts \(\Delta \bar{\kappa}_{\text{CD, Luft}}\) mit \hyperref[eq:arithmetisches_mittel]{Arithmetisches Mittel (\ref*{eq:arithmetisches_mittel})} und \hyperref[eq:fehler_mittelwert]{Fehler des Mittelwerts (\ref*{eq:fehler_mittelwert})}:

\eq{errCD_k_mittelwert}{
    \Delta {\overline{\kappa}_{\text{CD, Luft}}} = \frac{1}{n} \sqrt{\sum (\Delta \kappa_{\text{CD, Luft}})^2} \, .
}

Damit kommen wir auf das Ergebnis für den Mittelwert des Adiabatenkoeffizienten der Luft:

\res{\kappa_\text{Luft}}{1,374}{0,004}{}

\section{Adiabatenkoeffizient nach Rüchardt}
\label{sec:A2}

Nun sollen die Adiabatenkoeffizienten von Luft und Argon bestimmt werden. Dies passiert über die Periodendauer \(T\) der Schwingung. Es wurde die Zeit für 75 Schwinungen gemessen, um die Periodendauer \(T\) zu bestimmen. Fehler der Periodendauer \(\Delta T\) werden mit \hyperref[eq:fehler_proportional]{Fehler für $f = ax$ (\ref*{eq:fehler_proportional})} berechnet:

\eq{errR_T}{
    \Delta T = \sqrt{frac{(\Delta T_i)^2 + \sigma^2}{n}}
}

Dabei wurden $n = 5$ Messungen für Luft und Argon durchgeführt. $T_i$ beschreibt den $i$-ten Messwert der Periodendauer, $\sigma$ die Standardabweichung der Messwerte. 
Zusätzlich muss nach ceq{R_p}. Der Fehler bestimmt sich nach \hyperref[eq:gauss_fehlfortpflanzung]{Gauß'scher Fehlerfortpflanzung (\ref*{eq:gauss_fehlfortpflanzung})} zu:

\eq{errR_p}{
    \Delta p = \sqrt{\left(\frac{\partial p}{\partial m} \Delta m\right)^2 + \left(\frac{\partial p}{\partial A} \Delta A\right)^2 + \left(\frac{\partial p}{\partial g} \Delta g\right)^2} \, .
}